\documentclass{article}

\usepackage{import}
\usepackage{transparent}
\usepackage{xcolor}
\usepackage{tcolorbox}
\usepackage{amsmath}
\usepackage{upgreek}
\usepackage{enumitem}
\usepackage{amssymb}
\usepackage[margin=0.25in]{geometry}
\usepackage{listings}

%Commands definitions
\newcommand{\setbackgroundcolour}{\pagecolor[rgb]{0.19, 0.19, 0.19}}
\newcommand{\settextcolour}{\color[rgb]{0.87, 0.87, 0.87}}
\newcommand{\invertbackgroundtext}{\setbackgroundcolour\settextcolour}

%Command execution
%If this line is commented, then the appearance remaind as usual.
\invertbackgroundtext

\lstset{
    language=C,
    basicstyle=\ttfamily,
    keywordstyle=\color{cyan},
    commentstyle=\color{green!70!black},
    stringstyle=\color{red},
    showstringspaces=false,
    tabsize=4,
    breaklines=true,
    breakatwhitespace=false,
}



\author{Lukasz Kopyto}
\title{Wyklad 5.}

\begin{document}
\maketitle


\section*{Rownanie rekurencyjne liniowe.}

$$x_{n} + p_{k-1}x_{n-1} + \dots + p_{1}x_{n-k+1} + p_{0}x_{n-k} = f(n)$$

Przyklady:

$x_{n} = 2x_{n-1} \iff x_{n} - 2x_{n-1} = 0$

$x_{n} = x_{n-1} + x_{n-2} \iff x_{n} - x_{n-1} - x_{n-2} = 0$

$x_{n} = x_{n-1} + n^{2} \iff x_{n} - x_{n-1} = n^{2}$

Rownanie rekurencyjne liniowe jednorodne- wtedy $f(n) = 0$

Pokazemy najpierw jak rozwiazac rownanie jednorodne, a potem jak z tego wyciagnac rozwiazanie dla pewnych rownan niejednorodnych.

Mozemy wypisac k szczegolnych rozwiazan rownania jednorodnego,

$x_{n} + p_{k-1}x_{n-1} + \dots + p_{1}x_{n-k+1} + p_{0}x_{n-k} = 0$

Rownanie to wymaga k warunkow poczatkowych: $x_{0}, x_{1}, \dots, x_{k-1}$.
Jezeli je znamy to na podstawie rownania rekurencyjnego mozemy obliczyc $x_{k}, x_{k+1}, \dots$

Niech nasze szczegolne rozwiazanie ma warunki poczatkowe:

$a_{n}^{0} = (1,0,0,\dots,0)$

$a_{n}^{1} = (0,1,0,\dots,0)$

$a_{n}^{2} = (0,0,1,\dots,0)$

i tak dalej.

Wtedy rozwiazanie z warunkami poczatkowymi $x_{0}, x_{1}, \dots, x_{k-1}$ ma postac
$$x_{n} = x_{0}a_{n}^{0} + x_{1}a_{n}^{1} + \dots + x_{k-1}a_{n}^{k-1}$$
Innymi slowy ciagi $a_{n}^{0}, a_{n}^{1}, \dots, a_{n}^{k-1}$ stanowia baze przestrzeni liniowej wszystkich ciagow spelniajacych rekurencje

$x_{n} + p_{k-1}x_{n-1} + \dots + p_{1}x_{n-k+1} + p_{0}x_{n-k} = f(n)$

Wymiar tej przestrzeni to k.

Chcemy znalezc lepsza baze. Aby to zrobic, przedstawimy zaleznosc rekurencyjna za pomoca operatora przesuniecia E:
$$(x_{0}, x_{1}, x_{2}, x_{3}, \dots) \overset{E}{\longrightarrow} (x_{1}, x_{2}, x_{3}, \dots)$$

Piszemy: $E\langle x_{n}\rangle = \langle x_{n+1}\rangle$, czyli np.

$x_{n+k} + x_{n+k-1}p_{k-1} + \dots + x_{n}p_{0} = 0$

jest rowne

$E^{k}x_{n} + E^{k-1}x_{n}p_{k-1} + \dots + x_{n}p_{0} = 0$

czyli

$(E^{k} + E^{k-1} + \dots + E^{2} + E)<x_{n}> = <0>$

\subsection*{Dodatek, dokladniej o anihilatorach}



\end{document}
